\chapter{Markov processes and simulation}
\label{ch:markov-simulacion}

\section{Two ways to describe uncertainty}
A Markov chain propagates probabilities among states; a simulation constructs
realizations of events and estimates measures through repetition. The former
can yield a stationary distribution through equations; the latter supports
more detailed networks and rules, at the cost of sampling error. QSB maintains
separate normalized models and results for the two procedures
\cite{ross2019introduction,law2015simulation}.

\section{Finite chains and the row convention}
\(P_{ij}\) is the probability of moving from state \(i\) to \(j\) in one
period. Each row sums to one. The distribution is a row vector:
\[
 s_{t+1}=s_tP,\qquad \pi P=\pi,\qquad \sum_i\pi_i=1.
\]
Transitions are assumed time homogeneous, and the next state depends on the
current one without additional information from the past. Costs \(c_i\) are
per state per period; their expectation is \(s_tc\). They are not transition
costs.

\subsection{Example: machine availability}
\path{manual/examples/prob_markov.json} defines Operativa (operational) and
Parada (stopped), costs \((0,100)\), initial distribution \((1,0)\), a
horizon of three, and
\[
 P=\begin{pmatrix}0.9&0.1\\0.4&0.6\end{pmatrix}.
\]
From the repository root:
\begin{lstlisting}
.build/debug/qsb validate-markov-json manual/examples/prob_markov.json
.build/debug/qsb solve-markov-json manual/examples/prob_markov.json
\end{lstlisting}
The results are:
\begin{center}
\begin{tabular}{rrrr}
\toprule
Period & Operativa & Parada & Expected cost \\
\midrule
0 & 1 & 0 & 0 \\
1 & 0.9 & 0.1 & 10 \\
2 & 0.85 & 0.15 & 15 \\
3 & 0.825 & 0.175 & 17.5 \\
\bottomrule
\end{tabular}
\end{center}
The stationary distribution is \((0.8,0.2)\), with expected cost \(20\)
per period. This is a long-run mean, rather than a cumulative cost over three
periods. This chain is irreducible and aperiodic, so its transient distributions
converge to it. In general, stationary uniqueness and convergence from every
initial distribution are not equivalent: a periodic chain may oscillate.

QSB solves a stationary linear system and propagates matrices; it reports
\texttt{exact} as the algorithmic character, using floating-point arithmetic.
The method requires a unique stationary distribution and returns an error if
its system is singular. It does not choose an arbitrary mixture for chains
with several closed classes. When no initial distribution is supplied, it
does not invent one: transient output is empty. Validating rows and dimensions
does not guarantee that the stationary system is solvable.

In macOS, the editor supports names, costs, optional initial probabilities,
the transition matrix, and the horizon. Check row orientation before solving;
the result shows the stationary distribution, propagation, costs, and
assumptions. Editing states keeps dimensions consistent; invalid text remains
visible until corrected. JSON remains available for interchange.

\section{Discrete-event simulation}
A network contains entity sources, queues, servers, and gates. Routes include
probabilities and transfer times; service rules may depend on entity type.
QSB supports constant, exponential, normal, uniform, and triangular
distributions. Negative normal samples are clipped to zero; they are not
resampled. This clipping changes the mean and variability relative to an
unclipped normal distribution. Parameters represent distribution times or
values, rather than rates per time unit. For example, an exponential mean of
\(0.25\) hours corresponds to four arrivals/hour.

The educational backend models finite FIFO queues, alternative servers,
probabilistic routing, transfers, and assembly synchronization. Its outputs
are \texttt{approximate} estimates, even when a fixed seed makes them
reproducible. There is no warm-up deletion, interactive event trace, or
guarantee of steady state. No external simulation backend is available either.

\subsection{Example: a simulated service counter}
\path{prob_simulacion.json} connects source Cliente (customer), queue Espera
(waiting) with twenty places, and server Atención (service). Each arrival
brings one entity, interarrival time is exponential with mean \(0.25\)
hours, and mean exponential service time is \(1/6\) hours. Transfers take
zero time.
\begin{lstlisting}
.build/debug/qsb validate-simulation-json manual/examples/prob_simulacion.json
.build/debug/qsb solve-simulation-json manual/examples/prob_simulacion.json
.build/debug/qsb solve-simulation-json manual/examples/prob_simulacion.json --replications 3
\end{lstlisting}
The default run lasts \(1000\) hours, with seed \(1\), and no warm-up
removed. It generates \(4005\) entities, completes \(4004\), rejects
none, and estimates mean queue length \(1.202238\), maximum \(15\),
and utilization \(0.664470\). The uncompleted customer may be in the queue
or in service at termination: it is not necessarily rejected.

With three independently seeded replications, mean queue length is
\(1.229240\), with approximate normal \(95\%\) interval
\([1.146668,1.311813]\); utilization is \(0.659950\), with interval
\([0.638179,0.681722]\). Reported aggregate counters are rounded means:
\(4006\) generated and \(4005\) completed, rather than the sum of three
runs. Maxima, however, retain the maximum across replications. The base seed,
horizon, and number of replications appear in the output. The CLI exposes
replications but no options for choosing a seed or horizon in these commands;
those values come from the core's default options.

The interval formula uses \(\bar x\pm1.96s/\sqrt R\). With few
replications, the normal approximation may be poor; increasing \(R\) reduces
sampling error, but does not correct startup bias or an incorrect distribution.
Do not interpret the interval as a guarantee that the model describes reality.
The analytical M/M/1 in the queuing chapter gives \(L_q=4/3\) and
\(\rho=2/3\). The simulation has a finite queue and finite horizon:
closeness is useful for checking units, but does not require exact equality.
Three replications do not justify inferring a failure because an interval does
not contain the analytical value, either.

\section{macOS workflow}
In the simulation definition, complete \texttt{Components}, \texttt{Routes}, and \texttt{Service Rules}. Enter the parameters of each
distribution separately. Names must match route references and entity types.
Complete tables can be exchanged with spreadsheets; pasting checks each
section before replacing it. In the Run panel, select between one and one
hundred replications, validate, and run. The interface shows computation status
and prevents duplicate runs; changing the model invalidates a pending result.
This does not stop a calculation already started in the core.

The results panel shows entities generated/completed, completion rate, horizon,
seed, queue, and utilization, with intervals when available. Consult tables
and JSON as well as charts. Keep the network, rules, seed, horizon, and
replications when comparing scenarios: changing any of these may change the
interpretation of the difference.

\section{Exercises}
\begin{enumerate}
\item Remove the initial distribution from the Markov case and check that
only stationary analysis remains. Then change the matrix to the identity and
observe the singular-system error.
\item Double the simulated mean service time and compare completions,
rejections, and utilization. Explain why queue capacity is not equivalent to
service capacity.
\item Repeat the case with ten replications and compare interval widths.
Distinguish variation between runs from differences between model assumptions.
\end{enumerate}
